\section{Pre-Registered Strict Evaluation Protocol}
\label{app:protocol}

\paragraph{Frozen hypotheses.}
H1 tests whether temperature scaling lowers ECE on event-disjoint test disasters while leaving accuracy unchanged up to numerical ties.
H2 tests whether difference attention improves damage accuracy and calibration on event-disjoint disasters relative to concatenation under matched training.
H3 tests whether calibrated information-gain reinspection improves damage-decision accuracy per added sensing action relative to no reinspection.
H4 tests whether any gain remains relative to random and fixed-descent controls.
These hypotheses are evaluated separately; H1 does not imply H3.
Under the frozen checkpoints, H1 is supported on the test partition (ECE $0.095\rightarrow0.037$, accuracy unchanged) but not on holdout (ECE $0.046\rightarrow0.084$).
H2 is not supported as a uniform accuracy or holdout-calibration improvement; both fusions remain near $0.765/0.803$ accuracy.
H3--H4 are reported as underpowered rather than positive: the strict six-policy suite contains only one unique evidence-positive item with genuine triggers, observed $\Delta U=0$, and all pairwise SR comparisons remain non-significant after Holm correction.

\paragraph{Data partition.}
Disaster events, not crops, are the unit of partition.
Training, calibration, model-selection validation, and final testing use disjoint event sets.
All pre/post pairs and all crops belonging to one building remain together.
Before training, the pipeline records event lists, manifest hashes, image hashes, building identifiers, and class distributions.
Exact and near-duplicate checks run across partitions.

\paragraph{Model comparison.}
Concatenation and difference attention use the same encoder initialization policy, optimizer, augmentation, crop size, number of epochs, and training examples.
Multiple pre-declared random seeds are run for both variants.
Temperature is fitted independently for each seed using only the calibration split.
Architecture selection uses validation data; the final event-disjoint test remains unopened until selection is frozen.

\paragraph{Navigation comparison.}
Every policy receives the same instruction, start state, imagery, navigation backbone, perception checkpoint, maximum number of navigation steps, and random seed.
Policies differ only in reinspection action selection.
We report no reinspection, random reinspection, fixed descent, entropy threshold, and calibrated information gain.
Both the number of extra observations and altitude changes are reported; matched-budget secondary analyses cap each policy at the same reinspection allowance.

\paragraph{Primary analysis.}
The unit of inference is the episode.
We report paired differences with 95\% bootstrap confidence intervals and paired randomization-test \(p\)-values.
Strict SR and evidence-conditioned judgment accuracy are primary outcomes.
ECE, Brier score, NLL, semantic SR, NE, SPL, uncertainty reduction, trigger coverage, and cost-normalized gain are secondary outcomes.
All denominators are explicit.
If fewer than a pre-declared minimum number of episodes contain valid evidence and are eligible to trigger, H3--H4 are reported as underpowered rather than positive.

\paragraph{Failure accounting.}
We separately count no-evidence episodes, false evidence, wrong-target stops, timeout, budget exhaustion, altitude-floor termination, and unresolved reinspection at episode end.
Results are stratified by disaster event, target damage class, instruction difficulty, and initial target distance.
No subgroup is promoted to a headline result unless it was pre-specified.
